% AgriFlow — Proposal v13
% Compile: pdflatex AgriFlow_Proposal_v13.tex  (twice for refs)
\documentclass[11pt,a4paper]{article}

\usepackage[utf8]{inputenc}
\usepackage[T1]{fontenc}
\usepackage[a4paper,margin=2.3cm]{geometry}
\usepackage{booktabs}
\usepackage{longtable}
\usepackage{array}
\usepackage{enumitem}
\usepackage{xcolor}
\usepackage[hidelinks]{hyperref}
\usepackage{titlesec}

\definecolor{agrigreen}{RGB}{27,94,32}
\titleformat{\section}{\Large\bfseries\color{agrigreen}}{\thesection}{0.6em}{}
\titleformat{\subsection}{\large\bfseries}{\thesubsection}{0.5em}{}
\setlist{nosep,leftmargin=1.3em}
\newcommand{\rp}{Rp\,}

\title{\textbf{\color{agrigreen}AGRIFLOW} \\[2pt]
\large Platform Intelijen Ketahanan Pangan Sub-Nasional \\[4pt]
\normalsize \emph{Deteksi · Prediksi · Distribusi}}
\author{Tim AgriFlow --- Chelsea, Hilmi, Monika, Irpan}
\date{Proposal v13 --- PIDI DIGDAYA $\times$ Hackathon 2026, Bank Indonesia \\
Problem Statement \#2: Platform Matching Demand--Supply Antarwilayah}

\begin{document}
\maketitle
\thispagestyle{empty}

\begin{abstract}
\noindent
AgriFlow adalah platform intelijen ketahanan pangan yang mencocokkan kabupaten \emph{surplus}
dengan kabupaten \emph{defisit} pangan secara cerdas dan adil. Berbeda dari proposal awal yang
masih berupa rancangan, versi ini melaporkan \textbf{sistem yang sudah berjalan}: tiga fungsi
inti (Deteksi anomali harga, Prediksi harga, Distribusi/matching) berjalan di atas \textbf{data
BPS \& PIHPS Jawa Timur asli per-kabupaten}, dengan \textbf{364 tes otomatis lulus}, dashboard
\& chatbot WhatsApp yang \textbf{live}. Prinsip kami adalah \emph{honest engineering} --- setiap
klaim dapat diverifikasi terhadap kode dan data, dan keterbatasan dinyatakan apa adanya. Yang
membatasi cakupan saat ini bukan kemampuan sistem, melainkan \emph{ketersediaan data publik
per-kabupaten}; begitu data tersedia, pipeline yang sama langsung memprosesnya.
\end{abstract}

\section{Ringkasan Eksekutif}
Setiap tahun Indonesia kehilangan \rp 213--551 triliun pangan akibat \emph{food loss \& waste}
(4--5\% PDB), dan \textbf{40\% kerugian terjadi di distribusi, bukan produksi} (Bappenas 2021;
FAO 2019). Di satu kabupaten petani membuang cabai karena harga jatuh; di kabupaten sebelah
harga melonjak 100--200\% karena kelangkaan (PIHPS BI). Pemda kerap baru tahu krisis 2--3 minggu
kemudian.

AgriFlow memecahkan ini sebagai \textbf{alat bantu keputusan untuk pemerintah} (Bapanas/Pemda),
\emph{bukan} marketplace komersial. Inti platform adalah \emph{matching engine} 4-lapis yang
sadar masa-simpan, jarak jalan nyata, dan \textbf{keadilan untuk daerah tertinggal} --- dilengkapi
deteksi anomali harga dan peramalan harga, serta akses lewat dashboard dan WhatsApp (termasuk
Bahasa Jawa). AgriFlow dirancang sebagai \emph{operational layer} langsung untuk GPIPS 2026
Bank Indonesia.

\section{Masalah}
\textbf{Paradoks pangan.} Indonesia produsen pangan besar, tetapi distribusi antar-wilayah belum
efisien dan belum adil. Tabel berikut merangkum konteks yang mendasari AgriFlow.

\begin{center}
\begin{longtable}{p{9.5cm}p{5.5cm}}
\toprule
\textbf{Fakta} & \textbf{Sumber} \\
\midrule
\endhead
Food loss \& waste \rp 213--551 triliun/tahun (4--5\% PDB) & Bappenas Kajian FLW 2021 \\
40\% food loss di distribusi, bukan produksi & Bappenas 2021; FAO 2019 \\
Jatim: 4,7 juta petani, PDRB pertanian tertinggi nasional & BPS Jatim 2024 \\
Disparitas harga cabai antar-kab Jatim 100--200\% & PIHPS BI \\
Food loss bisa memberi makan 61--125 juta orang & Badan Pangan Nasional 2022 \\
BI meluncurkan GPIPS Februari 2026 & Bank Indonesia, 11 Feb 2026 \\
Target swasembada (Asta Cita \#2) & Perpres 12/2025 RPJMN \\
\bottomrule
\end{longtable}
\end{center}

\textbf{Paradoks aksesibilitas.} Solusi digital pangan yang ada umumnya menyasar petani muda
melek-smartphone. Realitas Indonesia berbeda: jutaan petani lansia, pengguna \emph{feature phone}
pedesaan, dan penyandang disabilitas. AgriFlow dirancang inklusif sejak awal --- akses lewat
kanal yang sudah ada di tangan pengguna (WhatsApp teks, dan Bahasa daerah), bukan afterthought.

\section{Solusi: Tiga Fungsi yang Sudah Berjalan}
AgriFlow memproses satu sumber data bersama (harga PIHPS harian; produksi/konsumsi/populasi BPS
per-kabupaten) menjadi tiga lapis nilai:

\begin{itemize}
  \item \textbf{Deteksi --- anomali harga.} Mengidentifikasi lonjakan/anjlok harga dari deret
        harga harian 2021--2025 memakai metode robust \emph{Seasonal-Hybrid ESD} (deseasonalize
        $\rightarrow$ statistik MAD pada residual). Memberi sinyal dini ke Pemda.
  \item \textbf{Prediksi --- harga 30 hari.} Peramalan harga ke depan dengan interval ketidakpastian
        (P10/P90). Saat ini berjalan pada \emph{baseline} statistik; varian foundation-model
        (TimesFM 2.0) sedang dikemas sebagai peningkatan akurasi.
  \item \textbf{Distribusi --- matching engine 4-lapis.} Mencocokkan kabupaten surplus dengan
        defisit melalui penyaringan keras (komoditas, jarak jalan nyata, masa-simpan, volume),
        penilaian multi-objektif 5 dimensi, dan alokasi terbobot-keadilan untuk kabupaten ber-IPM
        rendah. Metodologi lengkap di \emph{Dokumen Arsitektur} terpisah.
\end{itemize}

\textbf{Aksesibilitas.} Dashboard web (peta + tabel + grafik) dan \textbf{chatbot WhatsApp} yang
sudah live --- pengguna dapat menanyakan harga, mencari pembeli/pemasok, prediksi, dan anomali
lewat chat, dalam \textbf{Bahasa Indonesia} maupun \textbf{Bahasa Jawa}.

\subsection{Mengapa AgriFlow berbeda}
AgriFlow adalah \emph{government decision-support}, bukan marketplace --- berbeda dari eNAM
(India, lelang harga) atau MealConnect (AS, penyelamatan pangan), dan dari B2B agritech yang
gagal karena margin tipis (mis. TaniHub). Karena insentif kami adalah KPI ketahanan pangan +
stabilitas harga + keadilan, dimensi \emph{equity} antar-kabupaten menjadi tujuan utama.

\section{Bukti \& Validasi --- Bukan Sekadar Klaim}
\begin{itemize}
  \item \textbf{Data asli.} Engine berjalan pada data BPS/PIHPS Jawa Timur per-kabupaten ---
        produksi (BPS Jatim), konsumsi per-kapita (Susenas BPS; untuk beras bahkan
        \emph{per-kabupaten}, bukan asumsi nasional), populasi (BPS), harga harian PIHPS
        2021--2025. Lima komoditas inti tervalidasi nyata: beras, cabai merah \& rawit, bawang
        merah \& putih.
  \item \textbf{Teruji.} 364 tes otomatis lulus, lintas OS (CI GitHub Actions). Pipeline
        reproducible.
  \item \textbf{Live.} Dashboard (Vercel) + API \& chatbot WhatsApp (Hugging Face Spaces) berjalan.
  \item \textbf{Jujur \& terukur.} Nilai keadilan divalidasi empiris vs baseline (greedy/uniform/
        proporsional). Temuan kunci: saat pasokan langka, alokasi serakah menelantarkan kabupaten
        termiskin (Sampang 0\%, Bangkalan 20\%), sedangkan AgriFlow menjamin keduanya 100\%
        \emph{dengan biaya cakupan agregat nol}. Detail metodologi, evaluasi, dan 14 sitasi paper
        ada di Dokumen Arsitektur.
\end{itemize}

\section{GPIPS 2026 Alignment --- Fondasi Strategis}
11 Februari 2026 BI meluncurkan GPIPS (evolusi GNPIP) dengan penekanan penguatan pasokan
struktural via 3 pilar. AgriFlow dirancang sebagai \emph{operational layer} langsung untuk
ketiganya.

\begin{center}
\begin{longtable}{p{2.6cm}p{5.4cm}p{6.4cm}}
\toprule
\textbf{Pilar GPIPS} & \textbf{Mandat BI} & \textbf{Implementasi AgriFlow} \\
\midrule
\endhead
Pilar 1 --- Produksi & Peningkatan produksi, koordinasi pola tanam antar-wilayah & Deteksi
surplus/defisit per kabupaten; prediksi harga; dashboard koordinasi \\
Pilar 2 --- Distribusi & Efisiensi distribusi via Kerja Sama Antar Daerah (KAD) \& BUMN logistik &
Matching engine 4-lapis surplus $\leftrightarrow$ defisit; rekomendasi MOU antar-kab; smart routing \\
Pilar 3 --- Sinergi Kebijakan & Pemanfaatan data neraca pangan akurat Pusat--Daerah & Single
source of truth; export ke Bapanas/Bappeda/TPID; alert real-time \\
\bottomrule
\end{longtable}
\end{center}

\noindent\textbf{Strategic insight.} AgriFlow bukan vendor, melainkan mitra eksekusi GPIPS.
Sales cycle berubah dari ``beli produk kami'' menjadi ``kami menjalankan program yang sudah Anda
canangkan, dengan algoritma yang dapat diverifikasi terhadap data BPS dan kode terbuka.''

\subsection{Addressable Market (Pemerintah)}
\begin{center}
\begin{tabular}{lcr}
\toprule
\textbf{Segmen} & \textbf{Jumlah} & \textbf{Addressable Revenue/tahun} \\
\midrule
Kantor Perwakilan BI (KPw BI) & 46 & \rp 23--46 miliar \\
TPID Provinsi & 38 & \rp 7,6--19 miliar \\
TPID Kabupaten/Kota & 514 & \rp 38--77 miliar \\
Dinas Pertanian Kab/Kota & 514 & \rp 51--103 miliar \\
Bapanas Pusat & 1 & \rp 3--5 miliar \\
\midrule
\textbf{Total TAM Pemerintah} & --- & \textbf{\rp 120--250 miliar/tahun} \\
\bottomrule
\end{tabular}
\end{center}

\section{Keterbatasan \& Roadmap --- Yang Membatasi adalah Ketersediaan Data}
Penting digarisbawahi: \textbf{mesin AgriFlow sudah siap memproses data apa pun yang diberikan
kepadanya}. Yang membatasi cakupan saat ini adalah \emph{ketersediaan data publik
per-kabupaten}, bukan kemampuan sistem. Begitu sumber datanya terbuka, pipeline yang sama
langsung memprosesnya tanpa perubahan arsitektur.

\begin{center}
\begin{longtable}{p{6.2cm}p{8.6cm}}
\toprule
\textbf{Cakupan saat ini} & \textbf{Gerbangnya: ketersediaan data (bukan sistem)} \\
\midrule
\endhead
5 komoditas inti & Engine menerima komoditas apa pun. Komoditas lain menunggu data
\emph{produksi per-kabupaten} dipublikasikan BPS pada granularitas yang sama. \\
Tahun acuan 2022 & Tahun konsisten terbaru yang lengkap di semua sumber per-kabupaten. Tahun
lebih baru tinggal di-\emph{ingest} begitu BPS merilisnya. \\
Daging ayam \& telur belum & Data produksi broiler/ayam-ras per-kabupaten belum tersedia di sumber
publik; bukan keterbatasan engine. \\
Konsumsi cabai/bawang via angka nasional & Konsumsi \emph{per-kabupaten} untuk komoditas ini belum
dipublikasikan; konsumsi beras sudah per-kabupaten dan dipakai nyata. \\
Harga Tier-2 (kab non-IHK) & Panel Harga Bapanas sedang dalam pemeliharaan; saat data feed pulih,
30+ kabupaten tambahan langsung tercakup. \\
Skala provinsi (Jatim) & Engine siap; perluasan nasional 514 kab menunggu \emph{data feed
per-kabupaten nasional} terbuka, lalu optimasi skala (spatial partitioning). \\
\bottomrule
\end{longtable}
\end{center}

\noindent\textbf{Scaling up.} Peningkatan skala --- nasional, multi-komoditas penuh, real-time ---
\emph{dibatasi oleh laju keterbukaan data publik}, bukan oleh kesiapan teknis. Strategi kami:
buktikan nilai dulu pada skala provinsi dengan data nyata, lalu perluas seiring data tersedia.
Foundation-model forecasting (TimesFM 2.0) dan integrasi kanal suara/Bahasa daerah (Sahabat-AI,
IVR) adalah peningkatan terjadwal di fase berikutnya.

\section{Model Bisnis}
AgriFlow bukan proyek sosial yang bergantung hibah --- pendapatan dari 6 segmen via 8 aliran.

\begin{center}
\begin{tabular}{lcr}
\toprule
\textbf{Aliran Pendapatan} & \textbf{Target} & \textbf{Fase} \\
\midrule
1. Lisensi GPIPS Nasional & Bank Indonesia Pusat & Y2--Y3 \\
2. Langganan APBD Dinas & Dinas Pertanian + TPID & Y1--Y3 \\
3. B2B Data API (DaaS) & Bank, asuransi, FMCG & Y2--Y3 \\
4. White-Label Provinsi & Pemprov lain & Y2--Y3 \\
5. Custom Intelligence Reports & Think tank, Bapanas & Y1--Y3 \\
6. WhatsApp Premium Tier & Petani, pedagang & Y1--Y3 \\
7. Grants (Impact Capital) & BI Institute, UNDP, GIZ & Y1--Y3 \\
8. ESG \& Carbon Credit & Verra, Gold Standard & Y2--Y3 \\
\bottomrule
\end{tabular}
\end{center}

\subsection{Proyeksi 3 Tahun}
\begin{center}
\begin{tabular}{lrrr}
\toprule
& \textbf{Y1 (2026)} & \textbf{Y2 (2027)} & \textbf{Y3 (2028)} \\
\midrule
Total Revenue & \rp 650 jt & \rp 5.980 jt & \rp 16.100 jt \\
Net Income & \rp 450 jt & \rp 4.480 jt & \rp 11.100 jt \\
\bottomrule
\end{tabular}
\end{center}

\noindent Biaya operasional ramping: stack sengaja dipangkas (2 platform hosting; pgvector
alih-alih Qdrant; Gemini langsung alih-alih LangChain) --- \emph{pakai yang cukup, bukan yang
ramai}.

\section{Tim}
\begin{center}
\begin{tabular}{lll}
\toprule
\textbf{Nama} & \textbf{Peran} & \textbf{Fokus} \\
\midrule
Chelsea & Data Analyst & Forecasting (TimesFM), analisis harga \\
Hilmi & Data Architect & Arsitektur sistem, matching engine, validasi \\
Monika & UX Researcher & Riset pengguna, desain pengalaman, aksesibilitas \\
Irpan & Data Engineer & Pipeline data, backend, deployment \\
\bottomrule
\end{tabular}
\end{center}

\section{Penutup --- The Ask}
Dengan dukungan PIDI DIGDAYA $\times$ Hackathon 2026, AgriFlow siap menjalankan \emph{operational
layer} GPIPS 2026 BI: pilot 5 kabupaten Jawa Timur dengan matching engine, deteksi, dan prediksi
yang sudah berjalan di atas data nyata; menjangkau pengguna yang selama ini tertinggal lewat
WhatsApp dan Bahasa daerah; dan memperluas cakupan seiring data publik per-kabupaten terbuka.

\begin{center}
\emph{Deteksi · Prediksi · Distribusi --- untuk ketahanan pangan Indonesia.}
\end{center}

\end{document}
